\begin{answers}{ch:distributed}{分散学習：複数 GPU とクラスタで学習する}
\item 単純平均は $-3$ で、大バッチの $-4$ と一致しません。件数で重みを付けて平均すれば一致します。
$y=1$ の 1 件と、$y=3,5,7$ の 3 件に分けたとします。
装置 1 の勾配は $0-1=-1$、装置 2 の勾配は $\frac13(-3-5-7)=-5$ です。
単純平均は $\frac12(-1-5)=-3$ になり、1 件しかない装置 1 のデータを重く数えすぎています。
件数を重みにすると、$\frac{1\times(-1)+3\times(-5)}{4}=-4$ で、大バッチの勾配に戻ります。
式\eqref{eq:dist-avg}の変形は、どの装置も同じ $b=B/W$ 件を持つことを使っていました。
件数が違うときは、勾配の和（$-1$ と $-15$）と件数を別々に集め、全体の和 $-16$ を全体の件数 4 で割ります。
\item $W=8$ で 1.75 GB、$W=1024$ で約 1.998 GB です。
式\eqref{eq:dist-ring-bytes}の $2\frac{W-1}{W}S$ に $S=1$ GB を入れます。
$W=8$ では $2\times\frac78=1.75$、$W=1024$ では $2\times\frac{1023}{1024}\approx1.998$ です。
プロセスが 128 倍に増えても、1 ランクが送る量は 14\% しか増えず、2 GB を超えません。
\item 通信の時間は約 3 分の 1 になりました（8 ステップの短い実測）。
各パラメータの勾配を \code{torch.cat} で 1 本につなぎ、1 回の \code{dist.all_reduce} で送ってから、元の形に切り分けて戻しました。
4 プロセスでは 1 ステップあたり 31〜37 ms が 10〜14 ms に、2 プロセスでは約 6 ms が約 2 ms に減りました。
混んだ計算機で測ったので、どの値も本文より大きめです。
損失は、元の方法と小数第 6 位まで同じでした。
式\eqref{eq:dist-alpha-beta}の第 1 項（固定費）は、送る回数に比例します。
テンソル 20 個を別々に送ると 20 回分かかりますが、1 本にまとめれば 1 回分で済みます。
0 にならないのは、データ量に比例する第 2 項と、つなぐ・戻すための写しの時間が残るからです。
DDP のバケットは、この「まとめて送る」工夫を自動で行う仕組みです。
\item 約 2.2 GB です。
ZeRO-3 では、重み・勾配・オプティマイザの状態のすべてを $W$ 個に分けて持ちます。
1 パラメータあたり $2+2+12=16$ バイトなので、1 GPU の状態は $16N/W$ です。
$16\times70\times10^{9}\div512\approx2.19\times10^{9}$ バイトで、\file{zero_memory.txt} の最後の行（2.2 GB）と一致します。
分割なしの 1,120 GB、64 GPU の 17.5 GB と比べると、GPU を増やすほど 1 枚の負担が減ることが分かります。
ただし、これは活性化を除いた値です。
計算のときに一時的に集める層の重みの分も、別に必要です。
\item 63 個以上です。
式\eqref{eq:dist-bubble}に $P=8$ を入れると、バブルの割合は $\frac{7}{M+7}$ です。
これを $0.1$ 以下にするには $M+7\ge70$、つまり $M\ge63$ が必要です。
$M=63$ でちょうど 10\% になります。
一般に、バブルを 10\% 以下にするには、段数より 1 小さい数の 9 倍（$M\ge9(P-1)$）のマイクロバッチが要ります。
マイクロバッチを増やすと同時に持つ活性化が増えるので、1F1B のような順番の工夫と組み合わせます。
\end{answers}
